\notepage{77}{The linear Stark effect for hydrogen at n = 2}
Ignoring spin, the hydrogen $n=2$ subspace contains the four degenerate
states $\ket{200}$, $\ket{210}$, $\ket{211}$, and $\ket{21,-1}$.
Although every diagonal element
\[
 \mel{2\ell m}{eE_0Z}{2\ell m}=0
\]
vanishes by parity, the off-diagonal matrix elements need not vanish.
In the ordered basis
\[
 \{\ket{200},\ket{210},\ket{211},\ket{21,-1}\},
\]
the perturbation is
\[
 H'\big|_{n=2}=
 \begin{pmatrix}
 0&\Delta&0&0\\
 \Delta&0&0&0\\
 0&0&0&0\\
 0&0&0&0
 \end{pmatrix},\qquad
 \Delta=\mel{200}{eE_0Z}{210}.
\]
The phases of the basis vectors may be chosen so that $\Delta$ is real.
\smallheading{Exercise retained from the notes}
Calculate $\Delta$ explicitly using the hydrogen eigenfunctions. The
source suggests its textbook's Eq.~(13.1.27); the textbook and edition
are not specified on this page.
\smallheading{Diagonalization}
Use the basis
\[
 \left\{\frac{\ket{200}-\ket{210}}{\sqrt2},\quad
        \frac{\ket{200}+\ket{210}}{\sqrt2},\quad
        \ket{211},\quad\ket{21,-1}\right\}.
\]
The perturbation becomes
\[
 H'\big|_{n=2}=\diag(-\Delta,\Delta,0,0).
\]
The corresponding first-order shifts are $-\Delta$, $\Delta$, $0$, and
$0$. In this basis, off-diagonal perturbation elements between different
vectors of the original degenerate subspace vanish.
\editorial{As in note 74, the source's drawn diagonal matrix has the
signs interchanged relative to its stated order of the symmetric and
antisymmetric states. The matrix and basis order agree here.}
